%% ma: L12-B14-0135
%% lop: 12
%% chuong: 5
%% bai: 14
%% dang: 12.14.22
%% loai: TN4
%% muc: TH
%% dap_an: "C"
%% nguon: "Ảnh giáo viên gửi 25/11/2026 (tài liệu ôn tập Phương trình mặt phẳng), Phần I Câu 41"
%% kien_thuc: "Bài 14 (điều kiện hai mặt phẳng vuông góc)"
%% trang_thai: chinh-thuc
\begin{ex}
Trong không gian $Oxyz$, cho hai mặt phẳng $(\alpha)\colon m^2x-y+(m^2-2)z+2=0$ và $(\beta)\colon 2x+m^2y-2z+1=0$. Hai mặt phẳng $(\alpha)$ và $(\beta)$ vuông góc với nhau khi và chỉ khi
\choice{$|m|=\sqrt2$}{$|m|=1$}{\True $|m|=2$}{$|m|=\sqrt3$}
\loigiai{Hai vectơ pháp tuyến $\vec n_\alpha=(m^2;-1;m^2-2)$, $\vec n_\beta=(2;m^2;-2)$. $(\alpha)\perp(\beta)\Leftrightarrow\vec n_\alpha\cdot\vec n_\beta=0\Leftrightarrow2m^2-m^2-2(m^2-2)=0\Leftrightarrow m^2=4\Leftrightarrow|m|=2$.}
\end{ex}
